\section{Products of Addresses and Equality Certificates}\label{app:structure}
\subsection{Nearly Maximal Gaps in Polynomial Dimension}\label{app:productaddress}
For $\kappa\ge0$, an address block consists of $\kappa$ address bits and $2^\kappa$ data bits, and returns the addressed data bit. At $\kappa=0$ it is just one data bit. Let $f_j$ be address blocks on disjoint coordinates, with address lengths $\kappa_1,\ldots,\kappa_k$, and put
\[
 d=\prod_{j=1}^k f_j,\qquad
 b=\sum_{j=1}^k(\kappa_j+1),\qquad
 m=\sum_{j=1}^k(\kappa_j+2^{\kappa_j}),\qquad D=2^{\sum_j \kappa_j}.
\]
\begin{proposition}[Product address coefficient]\label{prop:productaddress}
For this contrast, $V_{\rm branch}(d,b)=1$ and $V_{\rm fixed}(d,b)=D=2^{b-k}$. Two rounds attain the branching coefficient.
\end{proposition}
\begin{proof}
Query all address bits first, and the $k$ selected data bits second. This evaluates $d$ within budget $b$, giving coefficient one by Jensen's inequality. For a fixed design, independently guess one data index in each block, query those $k$ bits and all address bits, and return $D d$ if every guess matches, zero otherwise. Completion has probability $1/D$ at every profile, so the conditional mean is $d$ and second moment is $D$.

For the converse let $\mu$ be uniform on the full cube and use $h=d$ in the variance dual. An arbitrary fixed form can allocate its queries among blocks in any way. If it observes $s_j$ address bits and $r_j$ data bits in block $j$, independence gives
\[
 \|\E_\mu[d\mid X_T]\|_2^2
 =\prod_{j=1}^k r_j2^{s_j-2\kappa_j}.
\]
If some $r_j=0$, this is zero. Otherwise $r_j\le2^{r_j-1}$ and $\sum_j(s_j+r_j)\le b$, whence
\[
 \prod_j r_j2^{s_j-2\kappa_j}
 \le 2^{\sum_j(s_j+r_j)-k-2\sum_j \kappa_j}
 \le 2^{-\sum_j \kappa_j}=D^{-1}.
\]
Since $\E_\mu[dh]=1$, the variance dual gives $V_{\rm fixed}\ge D$. This treats all fixed forms, including those that redistribute queries across blocks. The general minimax theorem then covers arbitrary biased estimators and history-dependent form selection.
\end{proof}

For any fixed $\kappa\ge1$, equal blocks give
\[
 m=\frac{\kappa+2^\kappa}{\kappa+1}b,\qquad
 \frac{V_{\rm fixed}}{V_{\rm branch}}=2^{\kappa b/(\kappa+1)}
 \quad\text{when }\kappa+1\text{ divides }b.
\]
In particular $m=3b/2$ permits $2^{b/2}$. More generally, for any fixed $\zeta>0$ one can choose a constant $\kappa$ so that the exponent is at least $(1-\zeta)b$ while $m=O_\zeta(b)$; the hidden dimension multiplier grows with $1/\zeta$. This does not solve the sharp frontier for a prescribed ratio $m/b$, such as $4/3$.

For every $b\ge2$, choose $\kappa=\lfloor\log_2 b\rfloor$, $k_0=\lfloor b/(\kappa+1)\rfloor$, and $\rho_0=b-k_0(\kappa+1)$. Use $k_0$ blocks of address length $\kappa$ and $\rho_0$ blocks of address length zero. Then $m=k_0(\kappa+2^\kappa)+\rho_0\le b^2$, and
\[
 \log_2(V_{\rm fixed}/V_{\rm branch})
 =k_0\kappa=b-k_0-\rho_0
 \ge b-\frac{b}{\kappa+1}-\kappa.
\]
Thus polynomial dimension already permits $2^{b-O(b/\log b)}$. The construction uses the standard product method; the content is the exact coefficient and the lower bound uniform over all query allocations.

\subsection{When the Full-Information Coefficient Is Attained}\label{app:equality}
Write $M=\|d\|_\infty>0$ and let $\mathcal H_\iota$ be the space of functions observable under interview $\iota$. Define the bounded-decomposition quantity
\[
 \mathcal C_\A(d)=\inf\left\{\sum_\iota\|g_\iota\|_\infty:
       g_\iota\in\mathcal H_\iota,\ \sum_\iota g_\iota=d\right\}.
\]
It is computed by a linear program: minimize $\sum_\iota u_\iota$ subject to $\sum_\iota g_\iota=d$ and $-u_\iota\le g_\iota(c)\le u_\iota$ on every cell. Its size is that of the enumerated profiles and menu, not polynomial in a general compact representation.

\begin{proposition}[Equality and a bounded-score certificate]\label{prop:equality}
For any identifiable $d$, $M\le \mathcal C_\A(d)$ and
\[
 M^2\le V(d,\A)\le \mathcal C_\A(d)^2.
\]
Consequently $\mathcal C_\A(d)=M$ suffices for $V(d,\A)=M^2$. Equivalently, it suffices that $d/M$ be a convex combination of observable functions with values in $[-1,1]$.
If $|d(x)|=M$ at every profile, then $V(d,\A)=M^2$ if and only if one interview in $\A$ computes $d$ exactly. In particular, a Boolean contrast on the full cube has $V_{\rm fixed}(d,b)=1$ if and only if it depends on at most $b$ coordinates.
\end{proposition}
\begin{proof}
The triangle inequality gives $\mathcal C_\A(d)\ge M$. For a decomposition of total norm $C$, choose $q_\iota=\|g_\iota\|_\infty/C$, dropping zero terms. Each returned score $g_\iota/q_\iota$ is bounded by $C$, giving $V\le C^2$. Jensen gives $V\ge M^2$.

For the last assertion the finite program attains its optimum: a bounded objective bounds every $g_\iota$, the simplex is compact, and the quadratic perspectives are closed at $q_\iota=0$. At an optimum with value $M^2$, for each profile,
\[
 d(x)^2=\left(\sum_\iota q_\iota H_\iota(x)\right)^2
 \le\sum_\iota q_\iota H_\iota(x)^2\le M^2.
\]
If $|d(x)|=M$ everywhere, equality in Jensen forces $H_\iota(x)=d(x)$ for every positive-weight interview and every profile. Any such interview computes $d$. The converse follows by using that interview alone. A fixed coordinate set computes $d$ precisely when it contains all of its relevant coordinates.
\end{proof}

For general real-valued $d$, equality in Jensen is forced only on the extremal set $E=\{x:|d(x)|=M\}$. Every used form must distinguish opposite-sign extremal profiles. That condition is necessary but does not by itself control second moments away from $E$. The bounded-decomposition test is a sufficient certificate; it is not asserted to be necessary for general $d$. Deterministic classifier differences take values in $\{-1,0,1\}$, so they are not covered by the constant-magnitude equivalence when zero contrasts occur.

Few extremal profiles alone do not imply a small gap. Take an even number $k\ge2$ of equal address blocks of length $\kappa\ge1$, let $f$ be their product, and put $\ell(x)=m^{-1}\sum_{i=1}^m x_i$. For $0<\epsilon<1$ define $d_\epsilon=(1-\epsilon)f+\epsilon\ell$. The only profile with $|d_\epsilon|=1$ is the all-positive profile: $\ell=\pm1$ only at the two constant profiles, and $f=1$ at both when $k$ is even. A mixture of the two-round exact score for $f$ and a uniformly sampled one-coordinate score for $\ell$ has magnitude one and conditional mean $d_\epsilon$, so $V_{\rm branch}=1$. Under the uniform law, $\E[f\ell]=0$ because any coordinate lies in only one block and another independent block has zero mean. The witness $h=f$ from Proposition~\ref{prop:productaddress} therefore gives
\[
 V_{\rm fixed}(d_\epsilon,b)\ge(1-\epsilon)^2 2^{k\kappa}.
\]
An exponentially large gap is therefore compatible with a non-Boolean contrast having a unique absolute maximum. Empirical near-equality should be established by certificates rather than inferred from extremal-profile counts.

\subsection{Optimizing a Root Augmentation}\label{app:weightedroot}
With $U=(X_{\rm root},Y)$ and $w=q_S^{-1}-1$, the variance identity in Appendix~\ref{app:augmentation} implies that, for fixed $q$, the population-optimal root function is
\[
 c^*(U)=\frac{\E[wt\mid U]}{\E[w\mid U]},
\]
with any value when the denominator is zero. This follows by minimizing the conditional quadratic $\E[w(t-c)^2\mid U]$. The weighted calibration counterpart uses the same form law as the original comparator and a single pseudo-observation at the calibration grand mean, weighted by the mean calibration weight. It remains globally unbiased. A reduction in its calibration criterion does not guarantee a reduction under the held-out law.
